% 04_footnotes_marginal.tex
% Footnotes and marginal notes for Digital Humanities
% Overleaf compatible: YES (standard packages)

\documentclass[12pt,a4paper]{article}
\usepackage[utf8]{inputenc}
\usepackage[T1]{fontenc}
\usepackage{lmodern}
\usepackage[english]{babel}
\usepackage{hyperref}
\usepackage{cleveref}
\usepackage{footmisc}
\usepackage{marginnote}
\usepackage{geometry}
\usepackage{fontspec}

% For marginal notes
\geometry{marginparwidth=3.5cm, marginparsep=0.5cm}

% Footnote configuration
\usepackage[stable]{footmisc}  % Stable footnotes for floats
\setlength{\footnotemargin}{1.5em}
\renewcommand{\thefootnote}{\arabic{footnote}}

\hypersetup{
    colorlinks=true,
    linkcolor=blue,
    citecolor=red,
    urlcolor=cyan
}

\title{Footnotes and Marginal Notes for Digital Humanities}
\author{Author Name}
\date{\today}

\begin{document}

\maketitle

\begin{abstract}
This example demonstrates footnote and marginal note techniques
essential for Digital Humanities scholarship: critical apparatus,
editorial notes, parallel texts, and manuscript annotations.
\end{abstract}

\section{Basic Footnotes}\label{sec:basic}
Standard footnote\footnote{This is a standard footnote.}.

Multiple footnotes\footnote{First footnote.}\footnote{Second footnote.}.

Footnote with citation\footnote{See \textcite{knuth1984texbook} for details.}.

Long footnote with multiple paragraphs\footnote{
This is a longer footnote that spans multiple paragraphs.

It can contain detailed commentary, source discussions,
or editorial remarks common in critical editions.
}.

\section{Symbolic Footnotes}\label{sec:symbolic}
\renewcommand{\thefootnote}{\fnsymbol{footnote}}
Symbol footnotes\footnote{First symbolic footnote.}
for front matter\footnote{Second symbolic footnote.}
or title pages\footnote{Third symbolic footnote.}.
\renewcommand{\thefootnote}{\arabic{footnote}}

\section{Critical Apparatus Footnotes}\label{sec:apparatus}
Apparatus footnotes for textual variants:

The manuscript reads ``manuscript variant''
\footnote{A: manuscript variant; B: printed edition variant; C: corrected variant.}
while the printed edition has ``printed variant''
\footnote{A: printed variant; B: manuscript variant.}.

Parallel passages\footnote{Parallel passage in MS Vat. lat. 1234, fol. 12v.}.

Conjectural emendation\footnote{Conjectured by Smith (2020); see \textcite{smith2023dh}.}.

\section{Marginal Notes}\label{sec:marginal}
Marginal notes appear in the margin\marginnote{This is a marginal note.}.

\marginnote{%
  \textbf{Editorial note:} This passage shows evidence of later interpolation.
  The hand changes at this point.
}[-1cm]% vertical offset

Inline marginal note\marginnote{Short note.}.

Longer marginal note with multiple lines\marginnote{%
This is a longer marginal note that can contain
editorial commentary, translation notes,
or cross-references to other witnesses.
}.

\section{Parallel Text / Facing Pages}\label{sec:parallel}
\begin{quote}
\marginnote{\textbf{Latin}:}\textit{Hic est textus latinus.}\\
\textbf{English}: This is the Latin text.\marginnote{Translation by J. Doe, 2023.}
\end{quote}

\section{Manuscript Annotations}\label{sec:manuscript}
Marginalia reproduction:
\begin{quotation}
\marginnote{\textit{Marginalia (Hand A):} ``Nota bene''}%
Main text with manuscript marginalia reproduced in the margin.
\end{quotation}

\section{Footnote in Minipage/Table}\label{sec:minipage}
Footnotes in tables require special handling:

\begin{table}[htbp]
\centering
\begin{minipage}{\linewidth}
\centering
\begin{tabular}{ll}
\hline
Column 1 & Column 2 \\
\hline
Data A\footnote{Table footnote A.} & Data B\footnote{Table footnote B.} \\
Data C & Data D \\
\hline
\end{tabular}
\caption{Table with footnotes}
\end{minipage}
\end{table}

\section{Endnotes Alternative}\label{sec:endnotes}
% For endnotes, use \usepackage{endnotes} and \let\footnote=\endnote
% Then \theendnotes at the end of document/chapter.

\end{document}